\RequirePackage{rotating}
\documentclass[12pt]{amsart}

\usepackage[margin=1.25in]{geometry}

\usepackage{amsmath}
\usepackage{amssymb}
\usepackage{amsthm,amscd}
\usepackage{url}
\usepackage{graphicx}
\usepackage{etoolbox}
\AtBeginEnvironment{quote}{\itshape}

\usepackage{enumitem}

\usepackage{tikz} 
\usetikzlibrary{automata,arrows,shapes,topaths,trees,backgrounds,decorations.pathreplacing,positioning}

\usepackage{hyperref}

\usepackage[all,cmtip]{xy}
\usepackage{varwidth}

%\newtheorem{theorem}{Theorem}[section]
\newtheorem{theorem}{Theorem}[section]
\newtheorem{lemma}[theorem]{Lemma}
\newtheorem*{lemstar}{Lemma}
\newtheorem{corollary}[theorem]{Corollary}
\newtheorem{conjecture}[theorem]{Conjecture}
\newtheorem{proposition}[theorem]{Proposition}
\newcounter{claims}[theorem]
\newtheorem{fact}[theorem]{Fact}

\theoremstyle{definition}
\newtheorem{definition}[theorem]{Definition}
\newtheorem{example}[theorem]{Example}
\newtheorem{exercise}[theorem]{Exercise}
\newtheorem*{openquestion}{Open Question}
\newtheorem{question}[theorem]{Question}
\newtheorem{notation}[theorem]{Notation}
\newcounter{cases}[theorem]
\newtheorem{case}[cases]{Case}
\newtheorem{subcase}{Subcase}
\numberwithin{subcase}{cases}

\theoremstyle{remark}
\newtheorem{remark}[theorem]{Remark}
\newtheorem{convention}[theorem]{Convention}
\newtheorem{claim}[claims]{Claim}
\newtheorem*{claimstar}{Claim}

\DeclareMathOperator{\res}{res}

\newcommand{\mc}[1]{\mathcal{#1}}

\newcommand{\Eres}{\upharpoonright\!\!\upharpoonright}
\newcommand{\la}{\langle}
\newcommand{\ra}{\rangle}
\newcommand{\conc}{\widehat{\ }}
\newcommand{\bft}{\mathbf{bf}}
\newcommand{\fneg}{\mathord{\sim}}

\newcommand{\treeres}[2]{\{\rotatebox[origin=c]{90}{$\rightsquigarrow$}_{#1}^{#2}\}}

\newcommand{\mbb}[1]{#1}

\newcommand{\A}{\mathsf{A}}
\newcommand{\B}{\mathsf{B}}
\newcommand{\C}{\mathsf{C}}
\DeclareMathOperator{\Frac}{Frac}
\newcommand{\D}{\mathsf{D}}

\newcommand{\Qp}{\mathbb{Z}_{(p)}}
\newcommand{\Q}[1]{\mathbb{Z}_{(#1)}}
\newcommand{\nBp}{\mathbb{Z}_{(p)}\mc{A}}
\newcommand{\Bp}{G^g_{\divs p}}
\newcommand{\tBp}{G^h_{\divs p}}
\newcommand{\Es}[2]{\mathbb{E}_{#2}^{#1}}
\newcommand{\Eq}{\mathrm{Equiv}}
\newcommand{\Znz}{\mathbb{Z}_{\neq 0}}
\newcommand{\Qnz}{\mathbb{Q}_{\neq 0}}
\DeclareMathOperator{\divs}{div}
\newcommand{\n}{\mathcal{N}}
\newcommand{\m}{\mathcal{M}}
\newcommand{\rE}{\mathrel{E}}
\DeclareMathOperator{\Mod}{Mod}

\DeclareMathOperator{\fractionfield}{Frac}

\newcommand{\f}{\mathsf{f}}

\newcommand{\Labels}{\mathtt{Labels}}


\DeclareMathOperator{\scope}{scope}
\DeclareMathOperator{\pred}{pred}
\DeclareMathOperator{\ran}{ran}
\DeclareMathOperator{\rank}{rank}
\DeclareMathOperator{\dom}{dom}
\DeclareMathOperator{\ext}{ext}
\DeclareMathOperator{\Th}{Th}
\DeclareMathOperator{\proves}{proves}
\DeclareMathOperator{\num}{num}
\DeclareMathOperator{\ind}{ind}
\DeclareMathOperator{\PC}{PC}
\DeclareMathOperator{\tup}{tup}
\DeclareMathOperator{\flen}{extlen}
\DeclareMathOperator{\Root}{root}
\DeclareMathOperator{\Cl}{Cl}

% Authors' macros::begin %%%%%%%%%%%%%%%%%%%%%%%%%%%%%%%%%%%%%%%%%%%%%%%
\newcommand{\cal}[1]{{\mathcal #1}}
\newcommand{\Sp}{\ensuremath{S}}
\newcommand{\SSp}{\ensuremath{\mathbb{S}}}
\newcommand{\Sfin}{\ensuremath{\mathbb{S}_{fin}}}
\newcommand{\fs}{\{0,1\}^*}
\newcommand{\N}{\mathbb{N}}
\newcommand{\NN}{\mathcal{N}}
%\newcommand{\lang}{\mathcal{L}}
\newcommand{\tuple}[1]{\langle #1 \rangle}
\newcommand{\lfs}{\mathsf{LFS}}
\newcommand{\lfsst}{\lfs^*}
\newcommand{\LFS}{\mathbf{LFS}}
\newcommand{\LFSst}{\mathbf{LFS}^*}
\newcommand{\cs}{\{0,1\}^\omega}
\newcommand{\zero}{\mathbf{0}}
\newcommand{\ptime}{\mathsf{P}}
\newcommand{\nptime}{\mathsf{NP}}
\newcommand{\exptime}{\mathsf{EXPTIME}}
\newcommand{\bpp}{\mathsf{BPP}}
\renewcommand{\L}{\mathcal{L}}
\renewcommand{\P}{\mathcal{P}}
\renewcommand{\H}{\mathcal{H}}
\newcommand{\Nc}{\mathcal{Nc}}
\newcommand{\comp}{\mathsf{COMP}}
\newcommand{\pspace}{\mathsf{PSPACE}}
\newcommand{\uh}{\upharpoonright}
\newcommand{\none}{\times}
\DeclareMathOperator{\tp}{tp}

\newcommand{\concat}[0]{\textrm{\^{}}}

\newcommand{\bigdoublewedge}{%
  \mathop{
    \mathchoice{\bigwedge\mkern-15mu\bigwedge}
               {\bigwedge\mkern-12.5mu\bigwedge}
               {\bigwedge\mkern-12.5mu\bigwedge}
               {\bigwedge\mkern-11mu\bigwedge}
    }
}

\newcommand{\bigdoublevee}{%
  \mathop{
    \mathchoice{\bigvee\mkern-15mu\bigvee}
               {\bigvee\mkern-12.5mu\bigvee}
               {\bigvee\mkern-12.5mu\bigvee}
               {\bigvee\mkern-11mu\bigvee}
    }
}

\renewcommand\labelenumi{(\theenumi)}

\begin{document}

\title{The Ritt problem is decidable}
\author{Jason Fang, James Freitag, Matthew Harrison-Trainor}
\address{Jason Fang, Department of Pure Mathematics, University of Waterloo, j7fang@uwaterloo.ca.}
\address{James Freitag, Department of Mathematics, Statistics, and Computer Science, University of Illinois at Chicago, jfreitag@uic.edu.}
\address{Matthew Harrison-Trainor, Department of Mathematics, Statistics, and Computer Science, University of Illinois at Chicago, mht@uic.edu.}

\thanks{Freitag was supported by NSF grants No. DMS-2452197 and DMS-2554149. Harrison-Trainor was supported by a Sloan Research Fellowship and NSF grant No. DMS-2452105. Harrison-Trainor received access to ChatGPT Pro through OpenAI’s academic researcher support program.}

\makeatletter
\def\blfootnote{\xdef\@thefnmark{}\@footnotetext}
\makeatother

\maketitle

\begin{abstract}
We give an affirmative solution to the Ritt problem. 
\end{abstract}

\subsubsection*{Declaration of AI Use}

The argument described in this paper was generated using ChatGPT-6 Astra. This manuscript was written by the authors and they take responsibility for the correctness of the results presented. As we describe in the paper, the argument follows a strategy based on work of Morrison \cite{Morrison79} that was pursued by Cohn \cite{Cohn1986Valuations}, of witnessing that one differential prime ideal is contained in another by finding a sequence of differential specializations from one to the other. A decision procedure for this condition is obtained from recursive axiomatizability, model completeness, and quantifier elimination developed by Scanlon \cite{Scanlon00} for the theory of nontrivially valued existentially closed monotone differential fields of equicharacteristic zero.

\section{Introduction}
\subsection{History and variations on the problems} 
In numerous works in analysis, algebraic differential equations, and the model theory of differential fields, the Ritt problem and its potential resolution have been considered. We mention several cases here, which we feel are particularly enlightening or significant from a mathematical or historical viewpoint. The reader who is unfamiliar with the theory of characteristic sets of differential ideals may find some of the technical remarks challenging without first consulting our short introduction (or reminder) about characteristic sets in Section \ref{charsetssection}. 

Throughout this section, we will consider systems of algebraic differential equations defined by the vanishing of some ordinary differential polynomials $f_1, \ldots, f_n$ perhaps in several variables over a differential field $K \subseteq \mc U$. By a \emph{singular solution} of the system, following Ritt \cite{ritt1936singular}, we will always mean an element $a \in \mc U$ at which the $f_i$ vanish and also for some $i$, the separant of $f_i$ vanishes at $a$. Understanding when such a singular solution is a specialization of the generic solution of the system is the lens through which Ritt approached this problem, though its origins are much earlier.

\begin{quote}
A number of the early continental analysts, notably Euler, Lagrange, Laplace, and Poisson, undertook the study of such singular solutions. They succeeded, by heuristic methods, in revealing certain interesting situations, but rigorous results, giving a complete description of the phenomena involved, were beyond their reach. \cite{ritt1936singular}
\end{quote}

A classical illustration is provided by the family
\[
F_{m,n}=(y')^m y''-y^n,\qquad m,n\geq 1,
\]
which Ritt attributes to Poisson \cite{Poisson1806} \cite[Chapter~III, \S38, p.~78]{Ritt50}. Working over $K=\mathbb{C}(x)$, the differential ideal 
\[
[F_{m,n}]:(y')^\infty
\]
is the prime differential ideal of the general component. Then $y=0$ is the only singular solution. While it is clear that $F_{m,n} \in [y]$, the Ritt problem demands that we understand when $[F_{m,n}]:(y')^\infty \subseteq [y].$

In this case it is possible via Ritt's work. As we mention below, Ritt solved the problem whenever the order of the equation is at most $2$. The low power theorem \cite[Chapter~III, \S20, p.~65]{Ritt50} implies
\[
[F_{m,n}]:(y')^\infty \subseteq [y] \quad\Longleftrightarrow\quad m<n.
\]
For instance, when $m=n=1,$ 
\[
(y'''-1)(3y'''-1)-y''y''''\in [F_{m,n}]:(y')^\infty,
\]
but this differential polynomial does not vanish at $y=0$. In terms described below, $0$ is not a specialization of the generic solution of the general component of $F_{m,n}=0.$

\subsection{From Laplace to Ritt}
In \cite{ritt1936singular}, Ritt describes the origins of the problem he considers coming from the work of Laplace \cite{Laplace1775}. Given a differential equation $F(y, y' ,\ldots y^{(n)})=0$ in a single unknown function, Laplace considered the two following problems: 
\begin{enumerate}
    \item determine whether an equation of lower order $A(y, y' ,\ldots y^{(r)})=0$ whose solutions satisfy  $F(y, y' ,\ldots y^{(n)})=0$ is \emph{contained in the general integral.}\footnote{The use of integral here is used in the way more modern writers would use solution.}
    \item to determine all the particular solutions of the given equation.\footnote{By a particular solution, Laplace meant a singular solution not contained in the general integral.}
\end{enumerate}
Ritt \cite{ritt1936singular} \cite{ritt1938great} discusses the meaning of Laplace's notion \emph{contained in the general integral}. If $P=[F]:H_F^\infty$ and $Q=[A]:H_A^\infty$, the hypothesis of (1) implies that $F \in Q$, but the first problem asks if $P \subseteq Q$. The second problem asks to find the prime components of the zero set of $F$. In \cite[49-52]{ritt1938great}, a more detailed discussion of the terminology and work of Poisson, Lagrange and Laplace is given. 

The problem is solved by Ritt in the case that $n=2$ in \cite{ritt1936singular}. Technically, Ritt worked over differential fields of meromorphic functions, but a later result of Seidenberg \cite{seidenberg1958abstract} would allow for a formal transfer of results that can be stated in the language of differential rings to arbitrary ordinary differential fields. Additionally, many of Ritt's results using the language of analysis were later put into purely algebraic terms by Kolchin and others. After summarizing the state of the art and the known special cases, Kolchin describes the progress on the general Ritt problem in his 1973 monograph \cite{Kolchin73}: 

\begin{quote}
    However, a solution of the general Ritt problem appears to be remote at present. 
\end{quote}


\subsection{Specialization and valuation} 
Starting around the 1970s with the work of Blum \cite{Blum1969, Blum1970}, Morrison \cite{Morrison78, Morrison79}, and Cohn \cite{Cohn77, Cohn1979, Cohn1986SecondOrder, Cohn1986Valuations} the Ritt problem was approached via valuation theory. This is fairly natural, since when $(L, \delta)$ is a differential field and $\mc O_L$ is a valuation ring with maximal ideal $\mathfrak m_L$ such that $\mc O_L$ is a differential ring and $\mathfrak m_L$ is a differential ideal, the residue map $\mc O_L \rightarrow k_L$ is a differential ring homomorphism. When $K$ is a differential subfield of $\mc O_L$, the residue map is injective on $K$ and so solution sets of differential equations with coefficients in $K$ are preserved by the residue map (here identifying $K$ with its image in the residue field). That is, if $\bar b = res (\bar a)$, then $\bar b$ is a differential specialization of $\bar a$ over $K.$ 

The adaptation from algebraic to differential algebraic geometry is not at all smooth with respect to valuation theory. For instance, in the classical situation, an integral domain $R$ is a valuation ring if and only if given $\alpha \in \Frac(R)$ either $\alpha$ or $\frac{1}{\alpha}$ is in the valuation ring. A consequence of this fact is that any algebraic specialization can be recovered as the residue map associated with a valuation, as we briefly explain. Given a specialization $\phi:R\to L$ of an integral domain $R$, let $\mathfrak p=\ker\phi$. There is a valuation $v$ of $\Frac(R)$ with the property that for $\alpha \in R,$ $v(\alpha) >0 $ for $\alpha \in \mathfrak p$ and $v(\alpha ) = 0 $ for $\alpha \in R \setminus \mathfrak p$. Then the homomorphism $\phi $ is given by the restriction of the residue map to the ring $R$. 

This picture begins to fall apart almost right away in the differential category. For instance, \cite[page 159]{Kolchin73} gives an example in which $R$ is a differential integral domain, $\phi: R \rightarrow F$ is a differential ring homomorphism to a differential field $F$, $L$ is a differential field containing $Frac(R)$, and $\eta \in L$ with the property that $\phi$ can not be extended to a differential homomorphism of $R\{\eta \} $ or $R\{ \eta^{-1} \}$ to any differential field containing $F$. It is this type of difficulty that necessitated the more complicated approach of Blum, Morrison, and Cohn whose technical details we deal with throughout sections \ref{sec:def-arrow} and \ref{specsec}.

In the above cited series of papers, Blum and Morrison developed the foundations of their valuative approach, and Cohn's series of papers applied their work to obtain partial results towards the Ritt problem. We continue this line of work in our solution, combining it with inputs described in the next several subsections as well as the model theory of valued differential fields, following Scanlon \cite{Scanlon00}, whose details we postpone discussion of until Section \ref{sec:def-arrow}. 

\subsection{Ideal containment, definability of irreducibility, and specialization}
Given a family of affine algebraic varieties $V \rightarrow W$ over an algebraically closed field $K$, the collection of $w \in W$ such that $V_w$ is irreducible is a definable subset of $W(K)$. The analogous fact for affine differential algebraic varieties over a differentially closed field $K$ can be seen to be equivalent to the Ritt problem.\footnote{This was perhaps first noticed in \cite{van1978model} and is proved in \cite{harrison2012nonstandard}.} The definability of irreducibility plays an important role in a number of works using the model theoretic approach. For instance, the positive resolution of the Ritt problem is expected to allow for more direct constructions of differential Chow varieties \cite{freitag2017differential}, new axiomatizations of the theory of differentially closed fields \cite{LeonSanchez16}, and computations involving ranks of differential equations \cite{freitag2025algebraic}. 

Even in considering the irreducibility problem, one is naturally led to the goal of studying specializations to the singular locus. For instance, echoing a formulation of Kolchin \cite{Kolchin73}, Hrushovski and Itai \cite{hrushovski2003model} write:

\begin{quote}
We do not know whether, in general, singularities included, irreducibility is a definable property in parameters. The question can be reduced to the following: given a point $p = (0,p_2, \ldots ,  p_m)$ and a Kolchin-closed set $X$, known to be irreducible on $x_1 \neq 0$, is $p$ in the closure of $X$? [sic]\footnote{We presume it was intended the final instance of $X$ should instead be $X \cap \{\bar x \, | \,  x_1 \neq 0 \}$ and this is the formulation we consider hence.} %$X \cap \{\bar x \, | \,  x_1 \neq 0 \}$? 
\end{quote}
%Since $X$ is already Kolchin-closed, read the last question as
%asking whether $p$ lies in the Kolchin closure of
%$U = X \cap \{\bar{x} \, | \, x_1 \neq 0\}$.
We explain this equivalence in detail, in part, because it reveals precisely where the potential difficulty in the Ritt problem arises in this particularly concrete setting.

To see that the above problem of Hrushovski and Itai can be solved given a solution to the Ritt problem, let $P$ be the ideal of differential polynomials over $K$ which vanish on $U = X \cap \{\bar x \, | \, x_1 \neq 0\}$; by hypothesis, this is a prime differential ideal.\footnote{Note that in principle we might be given this ideal in various different ways, but as detailed in Section \ref{effectivesection}, we can compute a characteristic set of the Kolchin closure of $U$ from a finite set of differential polynomials which cut out the closure of $U$ set theoretically.} Given the point $p$, let $Q$ be the ideal of differential polynomials over $K$ which vanish at $p$; the ideal of differential polynomials vanishing at a point is always prime. Then $p$ is in the Kolchin closure of $U$ if and only if $P \subseteq Q$. 

Conversely, let $P = [\mc C] : H_{\mc C} ^ \infty$ and $Q = [\mc D] : H_{\mc D} ^ \infty$. We first reduce away from two easier, specific, cases of the ideal membership problem. First, if some $c \in \mc C$ does not belong to $Q$, then $P \not \subseteq Q$. If every $c \in \mc C$ belongs to $Q$ and $H_{\mc{C}} \notin Q$, then $P \subseteq Q.$ To see this, note that for $f \in P,$ there is some $r$ so that $H_{\mc{C}}^r f \in [\mc C] \subseteq Q$. Then by the fact that $Q$ is prime, we have that $f $ must be in $Q$. 

Thus, the ideal containment problem is reduced to the case that $\mc C \subseteq Q$ and $H_{\mc{C}} \in Q$. Let $\bar b = (b_1 \ldots , b_n)$ be generic for $Q$ over $K$. Now, introduce a new first coordinate, call it $z$, and consider the point $p = (0, b_1, \ldots , b_n)$. Let \[X = \{(z,a) \, | \, \mc C(a) = 0, z = H_{\mc{C}}(a) \}, \quad U = X \cap \{(z,a) \, | \, z \neq 0\}.\]
$X$ is irreducible on $z \neq 0$ as $P = [\mc{C}] : H_{\mc{C}}^\infty$ is prime. Now it is easy to see that $P \subseteq Q$ if and only if $p \in \overline{U}^{Kol}$. 

\subsection{The computational approach}

The computational approach to the Ritt problem has, like all of the other approaches mentioned above, its origin in the characteristic-set methods of Ritt and Kolchin. The classical methods of Ritt and Kolchin compute, from a finite set $F$ of differential polynomials, a decomposition
\[
    \sqrt{[F]}=P_1\cap\cdots\cap P_r,
    \qquad
    P_i=[\mathcal C_i]:H_{\mathcal C_i}^{\infty},
\]
where each $P_i$ is prime and $\mathcal C_i$ is its characteristic set. In any existing implementation of the algorithm, the decomposition may contain redundant components, precisely because we were without a solution to the Ritt problem. (The Ritt-Kolchin algorithm requires the factorization of polynomials; in order to avoid this, a practical algorithm known as the Rosenfeld--Gr\"obner algorithm \cite{Ollivier90, BLOP95} computes a decomposition as above except that the $P_i=[\mathcal C_i]:H_{\mathcal C_i}^{\infty}$ need not be prime; such a decomposition is called regular. Such a decomposition is sufficient to test membership in $\sqrt{[F]}$. See \cite{BoulierLazardOllivierPetitot09} for some discussion of this.)

Golubitsky, Kondratieva, and Ovchinnikov
\cite[Theorem~1]{GolubitskyKondratievaOvchinnikov09} placed these questions in a common computability framework. Over a computable differential field of characteristic zero with a splitting algorithm,\footnote{This means that there is an algorithm for checking if an algebraic polynomial is irreducible.} they proved that the Ritt containment problem is equivalent to computing essential prime decompositions, testing primality of radical differential ideals given by finite generators, converting a prime's characteristic set into finite generators as a radical differential ideal, and testing whether a given polynomial is a zero divisor modulo a radical differential ideal. 

Their work clarified the precise computational obstruction and related it to several fundamental tasks in  computational differential algebra, but left the general decidability question unresolved. It is in their context that we give an algorithm to decide the Ritt problem.

\subsection{Main result}

\begin{theorem}
    Let $(K,\delta)$ be a computable ordinary differential field of characteristic zero with a splitting algorithm. The following problems are all decidable:
\begin{enumerate}
    \item Given a radical differential ideal specified by a finite set of generators, determine whether it is prime;
    \item Given a radical differential ideal specified by a finite set of generators, determine a non-redundant prime decomposition;
    \item Given a characteristic set of a prime differential ideal, find a finite set of generators for that ideal as a radical differential ideal;
    \item Given the characteristic sets of two prime differential ideals $P$ and $Q$, determine whether $P \subseteq Q$.
\end{enumerate}
\end{theorem}

We give an algorithm for (4), with the other decision problems being equivalent by the work of Golubitsky, Kondratieva, and Ovchinnikov \cite{GolubitskyKondratievaOvchinnikov09}. The algorithm is described in brief in Section \ref{sec:algorithm}, modulo three facts about specializations based on the work of Morrison \cite{Morrison78,Morrison79} and Scanlon \cite{Scanlon00}. These facts are proved in the subsequent sections.

\subsection{Limitations}

While we give an algorithm for deciding the Ritt problem, the algorithm is impractical for effective implementation in its current form without additional work. The algorithm exploits the fact that non-containment of prime ideals is easily seen to be semi-decidable (or computably enumerable), that is, given prime ideals $P$ and $Q$ (represented by their characteristic sets $\mc{C}$ and $\mc{D}$) there is an algorithm which will eventually halt if $P \nsubseteq Q$ and tell us so, but never halt if $P \subseteq Q$. This is simply a search for some $f \in P$ with $f \notin Q$, and we give no obvious bounds on how long this search might take. The other side of the algorithm shows that containment is also semi-decidable, with $P \subseteq Q$ being witnessed by a sequence of characteristic sets satisfying a relation $\mc{C} \rightsquigarrow \cdots \rightsquigarrow \mc{D}$. Again the algorithm searches through all possible sequences and there are no bounds on how long this might take or on how long this sequence might be in general.\footnote{When the Kolchin polynomials of the ideals being tested differ by $n \in N$, this number gives such a bound.} 

In particular, one step towards making the algorithm practical would be an effective version of Lemma \ref{lem:strict-step} for finding $P'$ strictly between $P$ and $Q$; the argument in this proof is ineffective, showing that we know such a $P'$ exists without computing it. Finally, we give no bounds on how long it might take to check $\mc{C} \rightsquigarrow \mc{D}$, though this should be in principle the easiest part of the argument in which to find bounds by investigating the quantifier elimination for existential formulas.

\subsection{Planned additions and future work}

We expect to soon add to this draft an argument for the existence of bounds version of the Ritt problem (and hence for the definability problem). We continue to investigate more efficient bounds for both the computational and existence-of-bounds/definability versions of the Ritt problem, as well as whether the solution can be extended beyond the ordinary case to the partial case.

\section{Background on differential algebra} \label{charsetssection}

We work over a differential field $K$ of characteristic zero with one derivation $\delta$. We denote the differential polynomial ring over $K$ in $n$ variables $x_1,\ldots,x_n$ by $K\{x_1,\ldots,x_n\}$. These are the polynomials with coefficients in $K$ in the variables $x_1,\ldots,x_n$ and their formal derivatives $\delta^\ell x_i$. Given a set of differential polynomials $F \subseteq K\{x_1,\ldots,x_n\}$, we denote by $(F)$, $[F]$, and $\{F\}$ respectively the ideal, differential ideal, and radical differential ideal generated. The degree of a differential polynomial $f \in K\{\bar{x}\}$ is the total degree of $f$ as a polynomial in $x$ and its derivatives. The order of a differential polynomial $f \in K\{\bar{x}\}$ is the maximum order of a differential operator appearing in $f$.

We rank the algebraic indeterminates $\delta^k x_i$ by putting $\delta^k x_i < \delta^\ell x_j$ if $(k,i) < (\ell,j)$ in the lexicographic order. Enumerating the algebraic indeterminates as $(Z_h : h \in \omega)$, we define the height of an algebraic indeterminate to be the $h$ such that that indeterminate is $Z_h$. Given a differential polynomial $f \in K\{\bar{x}\} - K$, the leader $v_f$ of $f$ is the highest-ranking algebraic indeterminate appearing in $f$, and the height of $f$ is the height of its leader. After writing $f$ as a polynomial in $v_f$, the leading coefficient is the initial $I_f$ of $f$. The leader of $I_f$ is of rank strictly less than $v_f$. The derivative of $f$ with respect to $v_f$ is called the separant $S_f$ of $f$.

We can extend the ranking on the algebraic indeterminates to a ranking on differential polynomials by the lexicographic order on the pair $(v_f,\deg_{v_f} f)$. Note that both the initial and separant of $f$ are of lower rank than $f$. We can also extend this ranking to a ranking of finite sequences of differential polynomials by putting $(f_1,\ldots,f_r) < (g_1,\ldots,g_s)$ if either (a) there is $k \leq \min(r,s)$ with $\rank(f_i) = \rank(g_i)$ for $i < k$ and $\rank(f_k) < \rank(g_k)$, or (b) if $r > s$ and $\rank(f_i) = \rank(g_i)$ for all $i \leq s$.

Given differential polynomials $f$ and $g$, we say that $f$ is reduced with respect to $g$ if no non-trivial derivative of $v_g$ appears in $f$ and the degree of $v_g$ in $f$ is strictly less than its degree in $g$. If $\rank(f) < \rank(g)$, then $f$ is reduced with respect to $g$, but not necessarily the other way around. A set of differential polynomials is autoreduced if each element is reduced with respect to the others. Autoreduced sets are always finite, and we can view them as ordered sequences by listing them in increasing height. This induces a ranking on the autoreduced sets of differential polynomials.

Given a set $S \subseteq K\{\bar{x}\}$, a characteristic set for $S$ is a lowest ranking autoreduced subset of $S$. Every proper differential ideal has a characteristic set. (We allow $\varnothing$ as a characteristic set for $S = (0)$ with $H_\varnothing = 1$.) Prime differential ideals are determined by their characteristic set by
\[ P = [\mc{C}] : H_\mc{C}^\infty  = \{ f \in K\{\bar{x}\} : \text{$H_{\mc{C}}^t f \in [\mc{C}]$ for some $t \in \mathbb{N}$} \}\]
where $H_{\mc{C}}$ is the product of the initials and separants of $\mc{C}$. There is a characterization of the characteristic sets of prime differential ideals. Since there is only one derivation, every autoreduced set is coherent.% autoreduced set $\mc{C}$ is coherent if whenever $f \neq g$ are in $\mc{C}$ are such that $\delta^k f = \delta^\ell g = Z_h$ and the height $h$ is minimal, $S_g \delta^k f - S_f \delta^\ell g \in (\mc{C})_h$ where $(\mc{C})_h$ is the ideal generated by $\mc{C}$ and the derivatives of $\mc{C}$ of height at most $h$.

\begin{lemma}[Lemma 2 of \S IV.9 of Kolchin \cite{Kolchin73}]\label{char-of-being-char-set}
    Let $\mc{C}$ be a finite subset of $K\{\bar{x}\}$. Then $\mc{C}$ is the characteristic set of a prime differential ideal if and only if
    \begin{enumerate}
        \item $\mc{C}$ is autoreduced, and
        \item $(\mc{C}) : H_\mc{C}^\infty$ is a prime ideal containing no non-zero elements that are reduced with respect to $\mc{C}$.
    \end{enumerate}
\end{lemma}

\section{Effective background} \label{effectivesection}

A computable differential field $(K,\delta)$ is a differential field whose domain is the natural numbers $\mathbb{N}$, and the addition, multiplication, and derivation operations on $K$ are computable (as functions on $\mathbb{N}$). One should think of the natural numbers as simply labeling the differential field elements.

We say that a field $K$ has a splitting algorithm if there is an algorithm to decide whether a polynomial $f(x) \in K[x]$ in one variable is irreducible. Since a splitting algorithm is required in order to decide whether an ideal $(f)$ generated by a single polynomial $f \in K[x]$ is prime, we ask the Ritt problem over a field with a splitting algorithm. The rational numbers and any computable algebraically closed field have splitting algorithms, and given a field $K$ with a splitting algorithm, classical results of Kronecker \cite{Kronecker82} (see also \cite{vanderWarden70}) give a splitting algorithm for an extension $K(a)$ from the splitting algorithm for $K$ together with the minimal polynomial of $a$ over $K$ (or alternatively from the information that $a$ is algebraically independent). Thus any finitely generated extension of $\mathbb{Q}$ has a splitting algorithm. In general, Rabin \cite{Rabin60} showed that any computable field can be embedded computably into its algebraic closure, and having a splitting algorithm is equivalent to the image of this map being computable.

\begin{question}[Ritt problem]
    Given a computable differential field $(K,\delta)$ with a splitting algorithm, is there an algorithm to determine whether a radical differential ideal $\{F\}$ specified by a finite set of generators $F \subseteq K\{x_1,\ldots,x_n\}$ is prime?
\end{question}

Golubitsky, Kondratieva, and Ovchinnikov \cite{GolubitskyKondratievaOvchinnikov09} show that the Ritt problem is equivalent to several other computational problems.

\begin{theorem}[Golubitsky, Kondratieva, and Ovchinnikov \cite{GolubitskyKondratievaOvchinnikov09}]
    The following problems are equivalent (in the sense that given an algorithm for one, we can find an algorithm for the others) over a computable differential field of characteristic zero with a splitting algorithm:
    \begin{enumerate}
        \item Given a radical differential ideal specified by a finite set of generators, determine whether it is prime;
        \item Given a radical differential ideal specified by a finite set of generators, determine a non-redundant prime decomposition;
        \item Given a characteristic set of a prime differential ideal, find a finite set of generators for that ideal as a radical differential ideal;
        \item Given the characteristic sets of two prime differential ideals $P$ and $Q$, determine whether $P \subseteq Q$.
    \end{enumerate}
\end{theorem}

In showing that the Ritt problem is decidable, we will give an algorithm for (4): Given the characteristic sets of two prime differential ideals $P$ and $Q$, there is an algorithm to decide whether $P \subseteq Q$. We will make use of several known algorithms over computable differential fields $K$ with a splitting algorithm:
\begin{enumerate}
    \item Given a prime differential ideal specified by its characteristic set $\mc{C}$, and a differential polynomial $f$, there is an algorithm to decide whether $f \in [\mc{C}] : H_{\mc{C}}^\infty$. This algorithm was first given by Ritt \cite[Chapter V §7]{Ritt50}.
    \item Given a finite set $\mc{C}$, using Lemma \ref{char-of-being-char-set} and standard algorithms from computational algebra we can determine whether $\mc{C}$ is the characteristic set of a prime differential ideal. Such an algorithm was first given by Ritt \cite[Chapter V \S 1]{Ritt50}.
    \item Given a radical differential ideal $\{ F \}$, the Ritt-Kolchin algorithm, which in the ordinary case is due to Ritt \cite[Chapter V \S 5]{Ritt50} (see also \cite{Kolchin73,Sit02,KondratievaLevinMikhalevPankratiev99}) produces a prime decomposition
    \[ \{ F \} = P_1 \cap \cdots \cap P_n\]
    where each prime differential ideal $P_i$ is represented by its characteristic set. This prime decomposition may be redundant. 
\end{enumerate}
The second and third algorithms require a splitting algorithm. For example, Golubitsky, Kondratieva, and Ovchinnikov \cite{GolubitskyKondratievaOvchinnikov09} remark that the Ritt-Kolchin algorithm requires one to check whether algebraic ideals $(\mc{C}):H_{\mc{C}}^\infty$, where $\mc{C}$ is a coherent autoreduced set, are prime and, if not, to find two polynomials not in the ideal whose product belongs to the ideal.

We will only use the first two algorithms, but we highlight the Ritt-Kolchin algorithm because of its role in tying together the equivalent characterizations of the Ritt problem: To test whether $\{F\}$ is prime, we can compute a prime decomposition, and then, if we could test containment of prime ideals (given their characteristic sets), we could find a minimal prime decomposition of $\{F\}$ and thus check whether it is prime.


\section{The algorithm (modulo three facts)}\label{sec:algorithm}

We briefly describe the high-level structure of the algorithm, black-boxing three facts for later. Given characteristic sets $\mc{C}$ and $\mc{D}$ for prime differential ideals $P$ and $Q$, we will define in Section \ref{sec:def-arrow} a relation $\mc{C} \rightsquigarrow \mc{D}$ which says, roughly, that ``there is a specialization of a non-singular $[\mc{C}]$-point (i.e., one at which $H_{\mc{C}}$ is non-zero) to a generic $Q = [\mc{D}]:H_{\mc{D}}^\infty$-point''.  This relation is dependent on the choice of characteristic set $\mc{C}$ for $P$ but not on the choice of the characteristic set $\mc{D}$ for $Q$. We will later prove three facts about this relation:
\begin{enumerate}
    \item If $\mc{C}$ and $\mc{D}$ are characteristic sets of prime differential ideals $P$ and $Q$ and $\mc{C} \rightsquigarrow \mc{D}$, then $P \subseteq Q$.
    \item Given $P \subseteq Q$ prime differential ideals with characteristic sets $\mc{C}$ and $\mc{D}$ respectively, either $\mc{C} \rightsquigarrow \mc{D}$ or there is a prime differential ideal $P'$ with characteristic set $\mc{C}'$ with $P \subsetneq P' \subseteq Q$ and $\mc{C} \rightsquigarrow \mc{C}'$.
    \item Given characteristic sets $\mc{C}$ and $\mc{D}$ for prime differential ideals, there is an algorithm to decide whether $\mc{C} \rightsquigarrow \mc{D}$.
\end{enumerate}
These are proved in Lemmas \ref{lem:prime-ideal-containment}, \ref{lem:strict-step}, and \ref{lem:decidability-of-arrow} respectively. We take these for granted for now. The idea is that given $P \subseteq Q$ with characteristic sets $\mc{C}$ and $\mc{D}$, (2) and the Ritt-Raudenbush basis theorem imply that there is a finite sequence of characteristic sets of prime differential ideals
\[ \mc{C} = \mc{C}_0 \rightsquigarrow \mc{C}_1 \rightsquigarrow \mc{C}_2 \rightsquigarrow \cdots \rightsquigarrow \mc{C}_n = \mc{D}. \tag{$*$}\label{cert}\]
Given any prime differential ideals $P$ and $Q$ specified by their characteristic sets $\mc{C}$ and $\mc{D}$, if there is such a sequence between $\mc{C}$ and $\mc{D}$, (1) allows us to conclude that $P \subseteq Q$. We call a sequence of characteristic sets $\mc{C} = \mc{C}_0,\mc{C}_1,\ldots,\mc{C}_{n-1},\mc{C}_n = \mc{D}$ for prime differential ideals $P = P_0,P_1,\ldots,P_n = Q$ satisfying \eqref{cert} a \textit{certificate} that $P \subseteq Q$.  (3) yields an algorithm to determine whether a particular sequence of characteristic sets of prime differential ideals is a certificate that $P \subseteq Q$.

We are given two prime differential ideals $P$ and $Q$ described by their characteristic sets $\mc{C}$ and $\mc{D}$ respectively, and must decide whether $P \subseteq Q$. On the one hand, we can compute a list of all characteristic sets of prime differential ideals---recalling that there is an algorithm to check whether a finite set is a characteristic set of a prime differential ideal---and search for a certificate that $P \subseteq Q$. For each sequence of characteristic sets we consider, we have an algorithm to check whether it is a certificate. On the other hand, we search for a polynomial $f \in K\{\bar{x}\}$ with $f \in P$ and $f \notin Q$, witnessing that $P \nsubseteq Q$. Membership in a prime differential ideal, represented by its characteristic set, is decidable. We can carry out these two searches simultaneously, and if $P \subseteq Q$ then the first search will terminate, while if $P \nsubseteq Q$ then the second search will terminate. Thus we can determine whether $P \subseteq Q$.

\section{Valued differential fields and the definition of $\rightsquigarrow$}\label{sec:def-arrow}

A valued differential field is a differential field equipped with a surjective valuation map $v:\mathcal K\to\Gamma \cup \{ \infty\}$ where $\Gamma$ is an ordered abelian group called the value group of $\mathcal K$. We use the additive notation for the valuation with $v(0) = \infty$. Write $\mathcal O_v=\{a\in\mathcal \mathcal K:v(a)\ge0\}$ for the valuation ring of $\mathcal K$ and $\mathfrak{m}_v=\{a\in\mathcal K:v(a)>0\}$ for the maximal ideal. Let $\kappa_v=\mathcal O_v/\mathfrak m_v$ be the residue field and $\res:\mathcal O_v\to\kappa_v$ is the quotient map. The valuation is \emph{nontrivial} if $\Gamma\ne\{0\}$.
A valuation has \emph{rank one} if its value group is nonzero and
order-isomorphic to an ordered subgroup of $(\mathbb R,+)$.
It has \emph{equicharacteristic zero} if both the valued field and its
residue field have characteristic zero. We say that $\mathcal K$ is \emph{monotone} (also sometimes called \emph{contractive}) if $v(\delta a) \geq v(a)$ for all $a \in \mathcal K$. In this case $\mathcal O_v$ is a differential ring and $\mathfrak m_v$ is a differential ideal, allowing us to equip the residue field $\kappa$ with the induced derivation making $\res$ a differential homomorphism.

Scanlon \cite{Scanlon00} introduces a theory $T$ of nontrivially
valued existentially closed monotone differential fields of
equicharacteristic zero. That paper describes the more general case of $D$-fields which, as described in Remark 2.2 of that paper, specializes to differential fields by taking $e = 0$ in the notation of that paper. We also take, as described just before the axiomatization in \S 5, the residue field to be a differentially closed field and the value group to be divisible. See also the paper \cite{Rideau19} of Rideau which discusses this specific case.

We work in the three-sorted language with sorts for the valued field $\mc{K}$, residue field $\kappa$, and value group $\Gamma$ equipped with the valuation $v : \mc{K} \to \Gamma\cup \{\infty\}$ and residue maps $\res : \mc{O}_{\mc{K}} \subseteq \mc{K} \to \kappa$. The derivation acts on both the valued field and the residue field. Scanlon \cite{Scanlon00} provides a recursive axiomatization of $T$ in \S 5:\footnote{Note that we can omit Scanlon's Axiom 5, $(\forall x \in K)\bigl[(\exists y \in K)\; y^n = x
\iff n \mid v(x)\bigr]$. Because the value group is divisible, this becomes equivalent to saying that whenever $x \neq 0$ there is $y$ with $y^n = x$. This follows from Hensel's lemma and the fact that the residue field is differentially closed. Given $x$, we can pick $c$ such that $v(c) = v(x) / n$, and let $u = x / c^n \in \mc{O}^\times$. Choose a non-zero $n$th root of $\res(u)$ in $\kappa$, and lift it to $a \in \mc{O}^\times$. Applying Hensel's lemma to $P(Y) = Y^n - u$, we have $v(P(a)) > 0$ and $v(P'(a)) = 0$. We find $b$ such that $b^n = u$, so that $(cb)^n = x$.}
\begin{enumerate}
    \item $(\mc{K},\delta)$ is a differential field of characteristic zero,
    \item $\Gamma$ is a non-zero divisible ordered abelian group,
    \item $(\kappa,\delta)$ is a differentially closed field of characteristic zero,
    \item $v : \mc{K} \to \Gamma \cup \{\infty\}$ is a valuation and $\res : \mc{O} \to \kappa$ is the residue map, both surjective,
    \item \textbf{Monotone:} for each $x \in \mc{K}$, $v(\delta x) \geq v(x)$, and for each $x \in \mc{O}$, $\delta \res(x) = \res(\delta x)$.
    \item \textbf{Enough constants}: for each $\gamma \in \Gamma$, there is $c \in \mc{K}^\times$ such that $\delta c = 0$ and $v(c) = \gamma$,
    \item \textbf{$\delta$-Henselian}: for every differential polynomial $P \in \mathcal O_v\{X\}$ and $a \in \mathcal O_v$,
whenever
\[
v(P(a))>0
\qquad\text{and}\qquad
v\!\left(
  \frac{\partial P}{\partial(\delta^i X)}(a)
\right)=0
\]
for some $i\geq 0$, there is $b\in\mathcal O_v$ such that
\[
P(b)=0
\qquad\text{and}\qquad
v(b-a)\geq v(P(a)).
\]
Here $\partial P/\partial(\delta^i X)$ is the formal
algebraic partial derivative, treating $X,\delta X,\delta^2 X,\ldots$
as independent indeterminates.
\end{enumerate}
Theorem 7.1 of \cite{Scanlon00} shows that $T$ is the model completion of the theory of  monotone valued differential fields of equicharacteristic zero, and that $T$ admits quantifier elimination. % trivially valued too if we need it

\begin{definition}
    Let $K$ be a differential field, which we equip with the trivial valuation. Let $\mc{D}$ be the characteristic set of a prime differential ideal $Q \subseteq K\{\bar{x}\}$. Let $b_1,\ldots,b_n$ be additional constant symbols in the residue field. Define
    \[ \Delta_{Q,K}(b)
 =\{f(b)=0:f\in Q\}
   \cup\{f(b)\ne0:f\in K\{\bar{x}\} - Q\} \]
   and
   \[
   T_{K,Q}
 =T\cup\operatorname{Diag}(K)\cup\Delta_{Q,K}(b).\]
\end{definition}
A model of $T_{K,Q}$ is an extension of $K$ to a model of $T$ with a generic $Q$-point in its residue field. Since the valuation on $K$ is trivial, $K$ is also a subfield of the residue field. We observe the following:
\begin{enumerate}
    \item $T_{K,Q}$ is consistent, as we may extend $K$ to $\Frac(K\{\bar{x}\}/Q)$ to include a generic point of $Q$ and, equipping it with the trivial valuation, embed it into a model of $T$ (using the fact that $T$ is the model completion of the theory of monotone valued differential fields of equicharacteristic zero);
    \item $T_{K,Q}$ is complete using the quantifier elimination of $T$, as quantifier-free statements about $b$ and $K$ in the residue field reduce to equations and inequations between differential polynomials decided by $\Delta_{Q,K}(b)$ and the diagram of $K$;
    \item $T_{K,Q}$ is recursively axiomatizable uniformly in the presentation of $K$ and the characteristic set $\mc{D}$ of $Q$. For this we use the fact that membership in $Q$ can be checked using the characteristic set $\mc{D}$ of $Q$.
\end{enumerate}

\begin{definition}
    Let $\mc{C}$ be the characteristic set of a prime differential ideal $P$. Let $b$ be a tuple of constant symbols in the residue field sort. Let
    \[ \Phi_{\mc{C},K}(b) = \exists x \in \mc{O} \; (\res(x) = b \text{ and } \bigwedge_{f \in \mc{C}} f(x) = 0\text{ and } H_{\mc{C}}(x) \neq 0)\]
    where $x$ is a tuple in the valued field sort.
\end{definition}

The sentence $\Phi_{\mc{C},K}(b)$ expresses that $b$ can be pulled back to a non-singular $P$-point (i.e., one at which $H_{\mc{C}}$ is non-zero) in the valued field sort.

\begin{definition}
    Let $\mc{C}$ and $\mc{D}$ be characteristic sets of prime differential ideals $P$ and $Q$ respectively. We define
$\mathcal C\rightsquigarrow_K\mathcal D$ if
\[
 T_{K,Q}\models\Phi_{\mathcal C,K}(b).
\]
\end{definition}

\begin{lemma}\label{lem:decidability-of-arrow}
    Let $K$ be a computable differential field. Given characteristic sets $\mc{C}$ and $\mc{D}$ of prime differential ideals $P$ and $Q$ of $K\{\bar{x}\}$, there is an algorithm to decide whether $\mc{C} \rightsquigarrow \mc{D}$.
\end{lemma}
\begin{proof}
    Since the theory $T_{K,Q}$ is complete and recursively axiomatizable, we can enumerate all of the consequences until we find that $T_{K,Q}\models\Phi_{\mathcal C,K}(b)$ or $T_{K,Q}\models \neg \Phi_{\mathcal C,K}(b)$.
\end{proof}

\begin{lemma}\label{lem:prime-ideal-containment}
    Let $K$ be a differential field and let $P$ and $Q$ be prime differential ideals of $K\{\bar{x}\}$ with characteristic sets $\mc{C}$ and $\mc{D}$ respectively. If $\mc{C} \rightsquigarrow \mc{D}$ then $P \subseteq Q$.
\end{lemma}
\begin{proof}
    Given $\mc{C} \rightsquigarrow \mc{D}$, we have that $T_{K,Q}\models\Phi_{\mathcal C,K}(b)$. Let $\mc{K}$ be a model of $T_{K,Q}$, and fix $a \in \mc{O}_{\mc{K}}$ witnessing $\Phi_{\mc{C},K}(b)$. Suppose that $f \in P$. Recalling that $P = [\mc{C}] : H_{\mc{C}}^\infty$, we can write $g = H_{\mc{C}}^k f$ for some $g \in [\mc{C}]$. Since $g(a) = 0$ and $H_{\mc{C}}(a) \neq 0$, we have $f(a) = 0$. Thus, taking residues (and recalling that the coefficients of $f$ are in $K$ and hence trivially valued) $f(b) = 0$. But $b$ is a generic $Q$-point, and so for $f \in K\{\bar{x}\}$, $f(b) = 0$ implies that $f \in Q$. We conclude that $P \subseteq Q$.
\end{proof}

\section{Constructing local specializations} \label{specsec}

In this section we will prove the key lemma.

\begin{lemma}\label{lem:strict-step}
Given characteristic sets $\mathcal C,\mathcal D$ of prime differential
ideals $P \subseteq Q$ respectively, either
$\mathcal C\rightsquigarrow_K\mathcal D$, or there is a prime differential
ideal $P'$ such that
\[
  P\subsetneq P'\subseteq Q
\]
and $\mathcal C\rightsquigarrow_K\mathcal C'$ for every characteristic set
$\mathcal C'$ of $P'$.
\end{lemma}

\begin{proof}
The first case to consider is the case when $H_{\mc{C}} \notin Q$. This is the simple case. Let $L = \fractionfield(K\{\bar{x}\}/Q)$ and let $a$ be a generic point of $Q$ in $L$. Putting the trivial valuation on $L$, we can embed $L$ in a model of $T$ (using the fact that $T$ is the model completion of the theory of monotone valued differential fields of equicharacteristic zero). Then as $P \subseteq Q$, $f(a) = 0$ for all $f \in P$, and $H_\mc{C}(a) \neq 0$ as $H_{\mc{C}} \notin Q$. The residue $b = \res(a)$ is also a generic point of $Q$. Since $T_{K,Q}$ is complete, to get $\mc{C} \rightsquigarrow \mc{D}$ it suffices to find, as we have, just one model of $T_{K,Q}$ satisfying $\Phi_{\mc{C},K}(b)$.

Now consider the case $H_{\mc{C}} \in Q$. We must construct an appropriate specialization witnessing that $\mc{C} \rightsquigarrow_K \mc{C}'$ for some characteristic set $\mc{C}'$ of a prime differential ideal $P'$. Consider the local ring
\[
  A=(K\{\bar{x}\}/P)_{Q/P},\qquad \mathfrak m=(Q/P)A,
\]
and let $a$ be $\bar{x} / P$ in $A$.
The derivation extends to this localization and preserves $\mathfrak m$. Note that as $H_\mc{C} \in Q$, $H_{\mc{C}}(a) \in \mathfrak{m}$ and as $H_\mc{C} \notin P$, $H_{\mc{C}}(a) \neq 0$. 

We make use of the following lemma which we will prove later in the section. 

\begin{lemma}\label{lem:local-specialization}
Let $(A,\mathfrak m)$ be a local integral domain containing $\mathbb{Q}$,
equipped with a derivation $\delta:A\to A$ such that
$\delta(\mathfrak m)\subseteq\mathfrak m$. For every
$0\ne h\in\mathfrak m$, there are a nontrivially valued monotone
differential field $(\mathcal K,\delta,v)$ and a differential
homomorphism $\eta:A\to\mathcal O_v$ such that
\[
 0\ne\eta(h)\in\mathfrak m_v,\qquad
 \ker(\res\circ\eta)\subseteq\mathfrak m.
\]
\end{lemma}

Applying Lemma \ref{lem:local-specialization} to $A$, we get a nontrivially valued monotone differential field $(L,\delta,v)$ and a differential homomorphism $\eta : A \to \mc{O}_v$ such that
\[ 0 \neq \eta(H_{\mc{C}}(a)) \in \mathfrak{m}_v \text{ and } \ker(\res \circ \eta) \subseteq \mathfrak{m}.\]
Note that as $K$ is a field, all non-zero elements of $K$ are units in $A$, and so $\eta |_K$ is an embedding of $K$ into $\mc{O}_v$, whence as the non-zero elements of $K$ are units in $\mc{O}_v$, the residue map yields an embedding of $K$ into the residue field $\kappa_v$. Let $\alpha = \eta(a)$, $\beta = \res(\alpha)$, and let $P'$ be the prime differential ideal
\[ P' = \{ f \in K\{\bar{x}\} : f(\beta) = 0\}.\]
For each $f \in [\mc{C}]$, $f(\alpha) = 0$ as $f(a) = 0$. Recall that $0 \neq \eta(H_{\mc{C}}(a)) = H_{\mc{C}}(\alpha) \in \mathfrak{m}_v$ and so $H_{\mc{C}}(\alpha) \neq 0$. Over $K$, $\beta$ is a generic point for $P'$.

Embed $(L,\mc{O}_v,\kappa_v)$ in a model
of $T$. The completeness of $T_{K,P'}$ together with the remarks above implies that $\mathcal C\rightsquigarrow_K\mathcal C'$ for every characteristic set
$\mathcal C'$ of $P'$. Lemma \ref{lem:prime-ideal-containment} gives $P \subseteq P'$. Since
$0\ne\eta(H_{\mc{C}}(a))\in\mathfrak m_v$, $H_{\mc{C}}(\beta) = 0$ whence $H_{\mc{C}}\in P' - P$, so the
inclusion is strict. If $f\in P'$, then
\[
  \res(\eta(f(a)))=f(\beta)=0.
\]
Thus $f(a)\in\ker(\res\circ\eta)\subseteq\mathfrak m = (Q/P) A$. The maximal ideal $(Q/P)A$ in the localization $A$ of $K\{\bar{x}\}/P$ contracts to the prime ideal $(Q/P)$. Thus $f\in Q$. We conclude that $P'\subseteq Q$.
\end{proof}

We now return to the proof of Lemma \ref{lem:local-specialization} which for the most part follows ideas of Morrison outlined in \cite{Morrison78} and \cite{Morrison79}.

We require the following definitions. Suppose that we have a differential field $L$ and a differential local subring $R \subseteq L$, which means that $R$ has a unique maximal ideal (which we denote $\mathfrak{m}_R$) and that this ideal is differential. The quotient map $R \to R/\mathfrak{m}_R$ gives a differential specialization. Another differential local subring $S \supseteq R$ dominates $R$ if its maximal ideal $\mathfrak{m}_S$ satisfies $\mathfrak{m}_S \cap R = \mathfrak{m}_R$, which means that the quotient map $S \to S/\mathfrak{m}_S$ extends $R \to R/\mathfrak{m}_R$. If the only differential subring of $L$ dominating $R$ is $R$ itself, then we say that $R$ is an $L$-maximal differential ring, and the map $R \to R/\mathfrak{m}_R$ is sometimes called a maximal differential specialization. By Zorn's lemma, every differential local subring of $L$ can be extended to an $L$-maximal differential ring.

\begin{proof}[Proof of Lemma \ref{lem:local-specialization}]
Let $(A,\mathfrak m)$ be a local integral domain containing $\mathbb{Q}$,
equipped with a derivation $\delta:A\to A$ such that
$\delta(\mathfrak m)\subseteq\mathfrak m$. Fix $0\ne h\in\mathfrak m$. We must construct a nontrivially valued monotone
differential field $(\mathcal K,\delta,v)$ and a differential
homomorphism $\eta:A\to\mathcal O_v$ such that
\begin{equation}
 0\ne\eta(h)\in\mathfrak m_v,\qquad
 \ker(\res\circ\eta)\subseteq\mathfrak m.
\end{equation}
Let $L=\Frac(A)$. The quotient $k=A/\mathfrak m$ is a
differential field and the quotient map
$\res_A:A\to k$, $a\mapsto a+\mathfrak m$, is a homomorphism of differential rings.

We begin by extending $A$ to an $L$-maximal differential ring  $R \subseteq L$ dominating $(A,\mathfrak{m})$, that is, with maximal ideal $\mathfrak{m}_R$ satisfying $\mathfrak{m}_R \cap A = \mathfrak{m}$.  When, as in our case, $R$ is an $L$-maximal differential ring and $L = \Frac(R)$, $R$ is sometimes called a differential valuation ring in the literature. As outlined in Section 4 of \cite{Morrison78} it is not necessarily the case that a differential valuation ring $R$ is a valuation ring. In Proposition 9 of \cite{Morrison78}, Morrison shows that for a differential valuation ring to be a (rank one) valuation ring suffices to have only one non-trivial differential prime ideal. Morrison \cite{Morrison79} also outlines a strategy that applies this to our situation.

We say that distinct differential prime ideals $\mathfrak{p}$ and $\mathfrak{q}$ are next neighbouring primes if $\mathfrak{p} \subseteq \mathfrak{q}$ and there are no differential prime ideals strictly in between. We apply the following proposition to obtain next neighbouring primes $\mathfrak{p}_- \subseteq \mathfrak{p}_+ \subseteq \mathfrak{m}_R$ of $R$ such that $h \in \mathfrak{p}_+$ and $h \notin \mathfrak{p}_-$.

\begin{proposition}[Morrison, Proposition 2.5 of \cite{Morrison79}]
    Let $R$ be a differential ring, and let $\mathfrak{p} \subsetneq \mathfrak{q}$ be prime differential ideals with $h \in \mathfrak{q}$ and $h \notin \mathfrak{p}$. Then there are next neighbouring prime differential ideals $\mathfrak{p}_-$ and $\mathfrak{p}_+$  with $\mathfrak{p} \subseteq \mathfrak{p}_- \subseteq \mathfrak{p}_+ \subseteq \mathfrak{q}$, $h \in \mathfrak{p}_+$, and $h \notin \mathfrak{p}_-$.
\end{proposition}

In our case, the argument is as follows. Choose a chain $\mathcal P$ of prime differential ideals of $R$,
maximal among chains containing $(0)$ and $\mathfrak m_R$ and lying
between them. Set
\[
 \mathfrak p_-=\bigcup_{\substack{\mathfrak p\in\mathcal P\\
                                  h\notin\mathfrak p}}\mathfrak p,
 \qquad
 \mathfrak p_+=\bigcap_{\substack{\mathfrak p\in\mathcal P\\
                                  h\in\mathfrak p}}\mathfrak p.
\]
Both collections are nonempty, since $h\ne0$ and
$h\in\mathfrak m_R$. The union and intersection are prime differential
ideals comparable with every member of $\mathcal P$, so maximality
puts them in $\mathcal P$. Moreover,
\begin{equation}\label{eq:adjacent-primes}
 (0)\subseteq\mathfrak p_-\subsetneq\mathfrak p_+
 \subseteq\mathfrak m_R,\qquad
 h\in\mathfrak p_+\setminus\mathfrak p_-.
\end{equation}
Every member of $\mathcal P$ is either contained in $\mathfrak p_-$
or contains $\mathfrak p_+$. A prime differential ideal strictly
between them could therefore be inserted into the chain, contrary
to maximality.

We now make use of the following two propositions of Morrison.

\begin{proposition}[Morrison, Proposition 2.1 of \cite{Morrison79}]
    Let $L$ be a differential field and let $R$ be an $L$-maximal differential ring. Let $\mathfrak{p}$ be a prime differential ideal of $R$. There is an $L$-maximal differential ring $A \subseteq L$ with maximal ideal $\mathfrak{p}$. $A$ consists of the elements $x$ such that $x \mathfrak{p} \subseteq \mathfrak{p}$.
\end{proposition}

\begin{proposition}[Morrison, Proposition 2.2 of \cite{Morrison79}]
    Let $L$ be a differential field and let $R$ be an $L$-maximal differential ring. Let $\mathfrak{p} \subseteq \mathfrak{q}$ be prime differential ideals of $R$ and let $A$ and $B$ be $L$-maximal differential rings with maximal ideals $\mathfrak{p}$ and $\mathfrak{q}$ respectively. Then
    \[ R \subseteq B \subseteq A.\]
\end{proposition}

Let $B_-$ and $B_+$ be the $L$-maximal differential rings with maximal ideals $\mathfrak p_-$ and $\mathfrak p_+$ respectively. We have
\[ R\subseteq B_+\subseteq B_-\subseteq L.\]
Let $V = B_+ / \mathfrak{p}_-$. Morrison shows that this is a rank one valuation ring, essentially because it has a unique non-zero prime differential ideal which reduces to a case studied in Morrison's previous paper \cite{Morrison78}.

\begin{proposition}[Morrison, Proposition 2.6 of \cite{Morrison79}]
    Let $R$ be an $L$-maximal differential ring, and let $\mathfrak{p}$ and $\mathfrak{q}$ be prime differential ideals of $R$ with $\mathfrak{p} \subseteq \mathfrak{q}$ such that $\mathfrak{p}$ and $\mathfrak{q}$ are next neighbouring primes. Let $A$ and $B$ be the $L$-maximal differential rings with maximal ideals $\mathfrak{p}$ and $\mathfrak{q}$ respectively. Then $B/\mathfrak{p}$ is a rank 1 valuation ring.
\end{proposition}

To spell out the maps, let
\[
 k_-=B_-/\mathfrak p_-,\qquad
 \rho_-:B_-\longrightarrow k_-,\quad b\longmapsto b+\mathfrak p_-,
 \qquad V=B_+ / \mathfrak{p}_- \cong \rho_-(B_+).
\]
Here $k_-$ is a differential field. The restriction
$\pi=\rho_-|_{B_+}:B_+\to V$ has kernel $\mathfrak p_-$, so
\[
 V\cong B_+/\mathfrak p_-,\qquad
 \mathfrak n=\rho_-(\mathfrak p_+)\cong
              \mathfrak p_+/\mathfrak p_-.
\]
Since $B_+$ is local, $\mathfrak n$ is the maximal ideal of $V$.
Set $\eta=\pi|_A$. All maps in the following diagram are differential,
and both squares commute:
\[
\begin{CD}
 A @>{\subseteq}>> B_+ @>{\subseteq}>> B_-\\
 @V{\eta}VV @VV{\pi}V @VV{\rho_-}V\\
 V @= V @>{\subseteq}>> k_- .
\end{CD}
\]
In particular, $\eta(a)=a+\mathfrak p_-$ for $a\in A$, with this
class regarded as an element of $V$. The image ring and its maximal
ideal satisfy
\[
 \delta V\subseteq V,\qquad
 \delta\mathfrak n\subseteq\mathfrak n,
\]
because $B_+$ and $\mathfrak p_+$ are preserved by $\delta$.

Let $\mathcal K=\Frac(V)$, with the induced derivation and the
valuation $v$ whose valuation ring is $V$. Thus
$\mathcal O_v=V$, $\mathfrak m_v=\mathfrak n$, and
$\kappa_v=V/\mathfrak n$. Equation~\eqref{eq:adjacent-primes} gives
\[
 0\ne\eta(h)\in\mathfrak n,
\]
so the valuation is nontrivial.
Every nonzero rational number is a unit of $A$, hence has unit
image in $V$ and nonzero residue. Consequently both $\mathcal K$
and $\kappa_v$ have characteristic zero. Given $a \in A$, $\eta(a)$ has zero residue precisely when
$a+\mathfrak p_-\in\mathfrak p_+/\mathfrak p_-$, or equivalently
$a\in\mathfrak p_+$. Hence
\[
 \ker(\res\circ\eta)=\mathfrak p_+\cap A
 \subseteq\mathfrak m_R\cap A=\mathfrak m.
\]
It remains to show that the valuation is monotone. Here, the fact that the valuation obtained is rank one plays a role. We use the following lemma which is a combination of Theorem 1 of \cite{Cohn77} and one of the remarks preceding it. This statement also appears as Corollary 6.1.2. of \cite{AschenbrennervandenDriesvanderHoeven17}. Cohn ascribes the essential ideas to Morrison \cite{Morrison78}.

\begin{lemma}[Cohn \cite{Cohn77}]\label{lem:rank-one-monotonicity}
Let $(E,\delta,v)$ be an equicharacteristic-zero valued differential
field whose value group has rank one. Suppose that its valuation
ring $V$ and maximal ideal $\mathfrak n$ satisfy
\[
 \delta V\subseteq V,\qquad \delta\mathfrak n\subseteq\mathfrak n.
\]
Then $v(\delta a)\ge v(a)$ for every $a\in E$.
\end{lemma}

\begin{proof}
Fix $0\ne a\in\mathfrak n$. Since
$\delta\mathfrak n\subseteq\mathfrak n$, Theorem 1 of \cite{Cohn77} (or, in our notation, Lemma~6.1.1 of \cite{AschenbrennervandenDriesvanderHoeven17}) gives,
for every nonnegative integer $r$,
\[
 (r+1)v(\delta a)\ge r v(a).
\]
If $\delta a=0$, the required inequality is immediate.
Otherwise, view the rank-one value group as an ordered subgroup of
$\mathbb R$. Dividing the displayed inequality in $\mathbb R$ by
$r+1$ and letting $r$ tend to infinity through the nonnegative
integers gives $v(\delta a)\ge v(a)$.
For a unit $a\in V^\times$, preservation of $V$ gives the inequality.
For $a\notin V$, apply the result to $a^{-1}\in\mathfrak n$ and use
$\delta(a^{-1})=-a^{-2}\delta a$. The case $a=0$ is immediate.
\end{proof}

Returning to the proof of Lemma~\ref{lem:local-specialization},
apply Lemma~\ref{lem:rank-one-monotonicity} to $E=\mathcal K$.
Thus $(\mathcal K,\delta,v)$ is monotone.
\end{proof}

% Finally, Morrison shows how to find neighbouring primes.

% \begin{lemma}
%     Let $R$ be a differential ring, and let $\mathfrak{p} \subsetneq \mathfrak{q}$ be prime differential ideals with $h \in \mathfrak{q} - \mathfrak{p}$. Then there are next neighbouring prime differential ideals $\mathfrak{p}^*$ and $\mathfrak{q}^*$  with $\mathfrak{p} \subseteq \mathfrak{p}^* \subseteq \mathfrak{q}^* \subseteq \mathfrak{q}$, $h \in \mathfrak{q}^*$, and $h \notin \mathfrak{p}^*$.
% \end{lemma}

% \begin{proof}[Proof of Lemma \ref{lem:local-specialization}]
% Let $(A,\mathfrak m)$ be a local integral domain containing $\mathbb{Q}$,
% equipped with a derivation $\delta:A\to A$ such that
% $\delta(\mathfrak m)\subseteq\mathfrak m$. Fix $0\ne h\in\mathfrak m$. We must construct a nontrivially valued monotone
% differential field $(\mathcal K,\delta,v)$ and a differential
% homomorphism $\eta:A\to\mathcal O_v$ such that
% \begin{equation}
%  0\ne\eta(h)\in\mathfrak m_v,\qquad
%  \ker(\res\circ\eta)\subseteq\mathfrak m.
% \end{equation}
% Let $L=\Frac(A)$. The quotient $k=A/\mathfrak m$ is a
% differential field and the quotient map
% $\res_A:A\to k$, $a\mapsto a+\mathfrak m$, is a homomorphism of differential rings.



% \smallskip\noindent
% \emph{Constructing a maximal extension of $\res_A$.}
% Consider pairs $(D,\mathfrak p)$, where $A\subseteq D\subseteq L$
% is a differential subring and $\mathfrak p$ is a prime differential
% ideal of $D$ with $\mathfrak p\cap A=\mathfrak m$.
% Order these pairs by
% \[
%  (D,\mathfrak p)\le(D',\mathfrak p')
%  \quad\Longleftrightarrow\quad
%  D\subseteq D'\ \text{ and }\ \mathfrak p'\cap D=\mathfrak p.
% \]
% The pair $(A,\mathfrak m)$ belongs to this set. Unions of chains of rings in this set are in this set. 
% Apply Zorn's lemma to get a maximal pair $(B,\mathfrak m_B)$.

% The localization $B_{\mathfrak m_B}$, obtained by inverting all
% elements of $B\setminus\mathfrak m_B$, is a differential subring of
% $L$. Its prime differential ideal
% $\mathfrak m_B B_{\mathfrak m_B}$ contracts to $\mathfrak m_B$ on $B$.
% It therefore gives a pair above $(B,\mathfrak m_B)$.
% Maximality forces $B=B_{\mathfrak m_B}$, so $B$ is local and
% $\mathfrak m_B$ is its maximal ideal. In particular,
% $k_B=B/\mathfrak m_B$ is a differential field. The quotient map
% $\res_B:B\to k_B$ fits into the commuting diagram
% \[
% \begin{CD}
%  A @>{\text{inclusion}}>> B\\
%  @V{\res_A}VV @VV{\res_B}V\\
%  k @>{\iota}>> k_B ,
% \end{CD}
% \]
% where $\iota(a+\mathfrak m)=a+\mathfrak m_B$ is an embedding because
% $\mathfrak m_B\cap A=\mathfrak m$.
% Any proper extension of $\res_B$ would, by taking its kernel, give a
% strictly larger pair in the above ordering. Hence $\res_B$ is
% $L$-maximal. We have proved explicitly that
% \[
%  A\subseteq B\subseteq L,\qquad
%  \mathfrak m_B\cap A=\mathfrak m.
% \]    
% This is what we will mean by extending the residue specialization to an
% $L$-maximal differential ring dominating $A$.

% \smallskip\noindent
% \emph{Choosing two adjacent prime differential ideals.}
% Two prime differential ideals
% $\mathfrak p_-\subsetneq\mathfrak p_+$ are called \emph{adjacent}
% if there is no prime differential ideal strictly between them.
% Choose a chain $\mathcal P$ of prime differential ideals of $B$,
% maximal among chains containing $(0)$ and $\mathfrak m_B$ and lying
% between them. Set
% \[
%  \mathfrak p_-=\bigcup_{\substack{\mathfrak p\in\mathcal P\\
%                                   h\notin\mathfrak p}}\mathfrak p,
%  \qquad
%  \mathfrak p_+=\bigcap_{\substack{\mathfrak p\in\mathcal P\\
%                                   h\in\mathfrak p}}\mathfrak p.
% \]
% Both collections are nonempty, since $h\ne0$ and
% $h\in\mathfrak m_B$. The union and intersection are prime differential
% ideals comparable with every member of $\mathcal P$, so maximality
% puts them in $\mathcal P$. Moreover,
% \begin{equation}\label{eq:adjacent-primes}
%  (0)\subseteq\mathfrak p_-\subsetneq\mathfrak p_+
%  \subseteq\mathfrak m_B,\qquad
%  h\in\mathfrak p_+\setminus\mathfrak p_-.
% \end{equation}
% Every member of $\mathcal P$ is either contained in $\mathfrak p_-$
% or contains $\mathfrak p_+$. A prime differential ideal strictly
% between them could therefore be inserted into the chain, contrary
% to maximality. This proves adjacency; this construction comes from
% \cite[Proposition~2.5]{Morrison79}.

% \smallskip\noindent
% \emph{Forming the valuation ring.}
% We now use the construction in
% \cite[Propositions~2.1--2.2 and~2.6]{Morrison79}. Applied to the
% $L$-maximal differential ring $B$ and the adjacent ideals just
% chosen, it supplies $L$-maximal differential subrings $B_+,B_-$
% of $L$ with
% \[
%  B\subseteq B_+\subseteq B_-\subseteq L,\qquad
%  \mathfrak m_{B_+}=\mathfrak p_+,\qquad
%  \mathfrak m_{B_-}=\mathfrak p_-,
% \]
% such that the image of $B_+$ in $B_-/\mathfrak p_-$ is an ordinary
% rank-one valuation ring of its fraction field.
% The displayed equalities of ideals are part of the result being
% used: $\mathfrak p_+$ itself is the maximal ideal of $B_+$, and
% $\mathfrak p_-$ itself is the maximal ideal of $B_-$, with all
% rings viewed inside $L$. Merely requiring their contractions to
% $B$ to be $\mathfrak p_+$ and $\mathfrak p_-$ would be a weaker
% statement. %Neither $B_+$ nor $B_-$ is being defined as a localization
% of $B$.

% To spell out the maps, let
% \[
%  k_-=B_-/\mathfrak p_-,\qquad
%  \rho_-:B_-\longrightarrow k_-,\quad b\longmapsto b+\mathfrak p_-,
%  \qquad V=\rho_-(B_+).
% \]
% Here $k_-$ is a differential field. The restriction
% $\pi=\rho_-|_{B_+}:B_+\to V$ has kernel $\mathfrak p_-$, so
% \[
%  V\cong B_+/\mathfrak p_-,\qquad
%  \mathfrak n=\rho_-(\mathfrak p_+)\cong
%               \mathfrak p_+/\mathfrak p_-.
% \]
% Since $B_+$ is local, $\mathfrak n$ is the maximal ideal of $V$.
% Set $\eta=\pi|_A$. All maps in the following diagram are differential,
% and both squares commute:
% \[
% \begin{CD}
%  A @>{\text{inclusion}}>> B_+ @>{\text{inclusion}}>> B_-\\
%  @V{\eta}VV @VV{\pi}V @VV{\rho_-}V\\
%  V @= V @>{\text{inclusion}}>> k_- .
% \end{CD}
% \]
% In particular, $\eta(a)=a+\mathfrak p_-$ for $a\in A$, with this
% class regarded as an element of $V$. The image ring and its maximal
% ideal satisfy
% \[
%  \delta V\subseteq V,\qquad
%  \delta\mathfrak n\subseteq\mathfrak n,
% \]
% because $B_+$ and $\mathfrak p_+$ are preserved by $\delta$.

% \smallskip\noindent
% \emph{Checking monotonicity and the required specialization.}
% Let $\mathcal K=\Frac(V)$, with the induced derivation and the
% valuation $v$ whose valuation ring is $V$. Thus
% $\mathcal O_v=V$, $\mathfrak m_v=\mathfrak n$, and
% $\kappa_v=V/\mathfrak n$. Equation~\eqref{eq:adjacent-primes} gives
% \[
%  0\ne\eta(h)\in\mathfrak n,
% \]
% so the valuation is nontrivial.
% Every nonzero rational number is a unit of $A$, hence has unit
% image in $V$ and nonzero residue. Consequently both $\mathcal K$
% and $\kappa_v$ have characteristic zero.

% We use the following lemma, which is the rank-one consequence of
% \cite[Lemma~6.1.1 and Corollary~6.1.2]{AschenbrennervandenDriesvanderHoeven17}. (really Cohn)

% \begin{lemma}[Rank-one monotonicity]\label{lem:rank-one-monotonicity}
% Let $(E,\delta,v)$ be an equicharacteristic-zero valued differential
% field whose value group has rank one. Suppose that its valuation
% ring $V$ and maximal ideal $\mathfrak n$ satisfy
% \[
%  \delta V\subseteq V,\qquad \delta\mathfrak n\subseteq\mathfrak n.
% \]
% Then $v(\delta a)\ge v(a)$ for every $a\in E$.
% \end{lemma}

% \begin{proof}[Proof of Lemma~\ref{lem:rank-one-monotonicity}]
% Fix $0\ne a\in\mathfrak n$. Since
% $\delta\mathfrak n\subseteq\mathfrak n$, \cite[Lemma~6.1.1]{ADH} gives,
% for every nonnegative integer $r$,
% \[
%  (r+1)v(\delta a)\ge r v(a).
% \]
% If $\delta a=0$, the required inequality is immediate.
% Otherwise, view the rank-one value group as an ordered subgroup of
% $\mathbb R$. Dividing the displayed inequality in $\mathbb R$ by
% $r+1$ and letting $r$ tend to infinity through the nonnegative
% integers gives $v(\delta a)\ge v(a)$.
% For a unit $a\in V^\times$, preservation of $V$ gives the inequality.
% For $a\notin V$, apply the result to $a^{-1}\in\mathfrak n$ and use
% $\delta(a^{-1})=-a^{-2}\delta a$. The case $a=0$ is immediate.
% \end{proof}

% Returning to the proof of Lemma~\ref{lem:local-specialization},
% apply Lemma~\ref{lem:rank-one-monotonicity} to $E=\mathcal K$.
% Thus $(\mathcal K,\delta,v)$ is monotone.

% Finally, $\eta(a)$ has zero residue precisely when
% $a+\mathfrak p_-\in\mathfrak p_+/\mathfrak p_-$, or equivalently
% $a\in\mathfrak p_+$. Hence
% \[
%  \ker(\res\circ\eta)=\mathfrak p_+\cap A
%  \subseteq\mathfrak m_B\cap A=\mathfrak m.
% \]
% This proves all the assertions of the lemma.
% \end{proof}



\bibliography{References}
\bibliographystyle{alpha}


\end{document}
